\section{Introduction}\label{sec:introduction}

A nef line bundle admits smooth metrics whose curvatures approach
semipositivity. A limit of those metrics may acquire singularities.
For an adjoint bundle, the question is whether the canonical and boundary
terms prevent logarithmic poles in a least singular semipositive metric.
We prove that they do for projective klt pairs, and use zero-Lelong
endpoint metrics to establish an injectivity theorem in ordinary
coherent cohomology.

We use the convention that a local metric weight $\phi$ gives squared
norm $|\xi|^2e^{-\phi}$. Semipositivity means that its curvature current
$i\partial\bar\partial\phi$ is nonnegative. Such a metric has
\emph{minimal singularities} if every other semipositive weight $\psi$
on the same line bundle satisfies $\psi\le\phi+C_\psi$ for a constant
$C_\psi$. Weights on rational line bundles are defined by clearing
denominators and dividing by the same positive integer. Lelong numbers
measure the logarithmic singularities of these weights; in particular,
zero Lelong numbers imply local integrability of $e^{-t\phi}$ for every
fixed $t>0$.

For a normal variety $H$ and an effective rational divisor $\Theta$,
the klt condition means that $K_H+\Theta$ is $\Q$-Cartier and, on a
log resolution $\pi:V\to H$, the crossing divisor $B$ in
$\pi^*(K_H+\Theta)=K_V+B$ has every coefficient less than one.
The actual rational line bundle, rather than its numerical class,
is the object in the following theorem.

\begin{theorem}\label{thm:adjoint-metric}
Let $(H,\Theta)$ be a normal connected projective complex klt pair, with
$\Theta$ an effective rational divisor.  Suppose that
$D_H=K_H+\Theta$ is a nef $\Q$-Cartier divisor.  For every projective log
resolution $\pi:V\to H$, the rational line bundle $N=\pi^*D_H$ has the following property: every semipositive singular Hermitian
metric with minimal singularities on $N$ has Lelong number zero
at every point of $V$. Such a metric exists. Neither $K_H$ nor $\Theta$ is required to be separately $\Q$-Cartier.
\end{theorem}

Consequently, for every minimal weight $\phi$ on $\pi^*(K_H+\Theta)$,
\[
                    \mathcal I(t\phi)=\cO_V\qquad(t>0),
\]
where $\mathcal I(t\phi)$ consists of holomorphic germs $f$ for which
$|f|^2e^{-t\phi}$ is locally integrable. We record the precise independence
of the metric and the Cartier multiple in
\Cref{cor:minimal-multiplier}.

Demailly--Peternell--Schneider constructed minimal-singularity metrics
on pseudoeffective line bundles and proved zero Lelong numbers in the
big and nef case \cite[Theorem~1.5 and Proposition~1.7]{DPS2001}.
Nefness alone does not give this conclusion. For example, on a ruled
surface over an elliptic curve, a nef divisor bundle can have its
divisor metric as a minimal metric, with a positive Lelong number
along the divisor \cite[Corollary~1.2 and Example~3.5]{Koike2015}.
Gongyo--Matsumura asked how to construct a zero-Lelong semipositive
metric on a nef log canonical bundle in their study of injectivity and abundance
\cite[Question~5.6]{GongyoMatsumura2017}.
\Cref{thm:adjoint-metric} gives a positive answer in the projective klt
setting on every projective log resolution, without bigness or
nonvanishing. Once one zero-Lelong semipositive metric exists, every
minimal metric is no more singular and therefore also has zero Lelong
numbers. Thus the universal assertion about minimal metrics expresses
the same existence conclusion on each pullback bundle.

\subsection{Interior boundaries and restriction of sections}

Our second result concerns a family of effective rational divisors
between two fixed endpoints. Their coefficients may cross integers.
The resulting changes in the integral part of the divisor give natural
inclusions of line bundles. Zero-Lelong metrics at the endpoints
control the induced map on ordinary $H^1$.

\begin{theorem}\label{thm:interior-injectivity}
Let $V$ be a smooth projective complex variety, let $L$ be an integral
divisor, and let $0\le C_0\le C_2$ be effective rational divisors whose
combined support is simple normal crossing.  Suppose that both
rational line bundles $L-C_0$ and $L-C_2$ admit singular Hermitian
metrics of semipositive curvature with zero Lelong numbers at every
point.  For a rational number $0<\lambda<1$, set
\[
 C_1=(1-\lambda)C_0+\lambda C_2,\qquad
 L_i=L-\floor{C_i}\quad(i=0,1).
\]
Then the natural inclusion of line bundles induces an injection
\[
 H^1\bigl(V,\cO_V(K_V+L_1)\bigr)
 \longrightarrow H^1\bigl(V,\cO_V(K_V+L_0)\bigr).
\]
\end{theorem}

The strict inequality $\lambda<1$ provides a positive coefficient in
the curvature comparison. Effectivity and the common crossing support
provide integrability for the fractional boundary weights. The theorem
concerns the displayed integral line bundles and the natural sheaf map;
its endpoint assumptions are on the actual rational bundles $L-C_i$.

One direct consequence explains the connection with restriction of
sections. Put $E=\lfloor C_1\rfloor-\lfloor C_0\rfloor$, an effective
integral divisor, which may be nonreduced. If $E\ne0$, then
\begin{equation}\label{eq:restriction-surjective}
 H^0\bigl(V,\cO_V(K_V+L_0)\bigr)
 \longrightarrow
 H^0\bigl(E,\cO_V(K_V+L_0)|_E\bigr)
 \quad\text{is surjective}.
\end{equation}
Indeed, the connecting homomorphism for
$0\to\cO_V(K_V+L_1)\to\cO_V(K_V+L_0)
\to\cO_V(K_V+L_0)|_E\to0$ has image equal to the kernel of
the injective $H^1$ map. Thus the comparison lifts sections from the
entire divisor scheme $E$ at once.

Classical root-cover injectivity, including a reduced boundary, appears
in Esnault--Viehweg \cite[Theorem~5.1]{EsnaultViehweg1992}.
Fujino's transcendental approach develops complete metrics on an
analytic complement and comparison with coherent cohomology
\cite[Section~3, Lemmas~3.1--3.2 and Claim~1]{FujinoInjectivity2012}.
Matsumura extends analytic injectivity to metrics with transcendental
singularities and obtains uniform estimates for primitives
\cite[Theorem~1.3 and Sections~5.2--5.3]{Matsumura2018}.
Our proof uses the same harmonic-form and local-primitive methods, with
a comparison of two endpoint weights. We prove the required uniform
primitive estimate and the return to ordinary cohomology explicitly.

\subsection{The two analytic arguments}

For the metric theorem, write $N=\pi^*(K_H+\Theta)=K_V+B$.
The adjoint hypothesis enters twice. The klt coefficients make the
adjoint density integrable, even where $B$ has negative exceptional
coefficients. Klt vanishing and nefness then give one ample bundle
$A_H$ on $H$ whose pullback $A=\pi^*A_H$ makes $mN+A$ globally
generated for every positive integer $m$ with $m(K_H+\Theta)$ Cartier.
Normalize a section $s$ by its integral
\[
 Q_m(s)=\int_V |s|^{2/m}e^{-a/m-b},
\]
where $a$ is a smooth weight of $A$ and $b$ the divisor weight of $B$.
Choose a normalized section maximizing its value at a proposed
positive-Lelong point.

The $L^{2/m}$ normalization has its origin in Narasimhan--Simha's
work on ample canonical bundles \cite{NarasimhanSimha1968}; its role
in the later envelope construction is explained in
\cite[Section~5]{BermanDemailly2012}.
Tsuji constructs canonical singular metrics using normalized sections
of $A+mK_X$ with a fixed sufficiently ample twist
\cite[Sections~1.2 and~2]{Tsuji2007}.
Berman--Demailly develop integral-normalized adjoint envelopes and
their algebraic approximation
\cite[Definition~5.3 and Proposition~5.19]{BermanDemailly2012}.
Their general pseudoeffective approximation allows twists $p_mA$ with
$p_m\to\infty$ and $p_m/m\to0$
\cite[Equation~(5.24)]{BermanDemailly2012}.
Here the fixed twist and the klt density support an extremal argument:
a localization estimate produces a better normalized section whenever
the minimal weight has a positive Lelong number.

The normalized extremal sections have a semipositive limit after the
twist is divided by $m$. Minimality bounds that limit by the chosen
minimal weight. Consequently a fixed sublevel set of the difference
of the weights has adjoint mass tending to zero. A localized
$\bar\partial$ equation, solved on an affine open set, produces a
holomorphic competitor whose weighted norm is bounded by that mass.
The exact localization constant is independent of $m$ and of the
affine open set. Positive Lelong number forces this competitor to have
the same value at the chosen point. A H\"older estimate whose final bound is independent of $m$ makes
its normalization strictly smaller, contradicting extremality.

Two details keep the construction uniform and valid on a log resolution.
The weak limits in the localization lemma are taken in fixed weighted
Hilbert spaces. Extension across negative exceptional coefficients is
performed on the normal base $H$ in the actual Cartier bundle
$m(K_H+\Theta)+A_H$, and then pulled back. The complete K\"ahler
$L^2$ estimate supplies the local analytic solution
\cite[Theorem~5.1]{Demailly2012}; the localization and extension
arguments are given in \Cref{sec:adjoint-metrics}.

For injectivity, regularize both endpoint metrics, retaining arbitrarily
small curvature loss and uniform integrability at every fixed exponent.
A logarithm of a sum of endpoint weights makes the inclusion of line
bundles a contraction. On the complement of the crossing divisor,
a complete K\"ahler metric with bounded local potentials allows
local $L^2$ solutions. Their holomorphic differences extend across the
divisor and define ordinary \v Cech classes. An open-mapping argument
on one fixed finite cover gives a uniform bound for primitives of
forms whose ordinary class vanishes.

A Bochner comparison then makes the image of a harmonic representative
small. All approximating representatives embed into a single fixed
Hilbert space; weak convergence there preserves the original ordinary
cohomology class. This is the final step from weighted estimates on the
complement to the coherent-cohomology injection. The argument is
independent of the adjoint metric theorem: it uses only the endpoint
metrics stated in \Cref{thm:interior-injectivity}.

\subsection{A fourfold application and organization}

The metric theorem also gives a short abundance argument for a
projective complex klt fourfold $(X,\Delta)$ with nef rational adjoint
and $\chi(X,\cO_X)\ne0$. Lazi\'c--Peternell's nonvanishing theorem
uses generalized algebraic singularities of a semipositive metric
\cite[Corollary~D]{LazicPeternell2020}; zero Lelong numbers satisfy
that condition with no divisorial part. With a first section obtained,
Gongyo--Matsumura's criterion gives semiampleness
\cite[Corollary~5.3]{GongyoMatsumura2017}.
The precise application is \Cref{cor:euler-abundance}; it has no
numerical-dimension restriction. For nef adjoints of numerical dimension at most one,
Liu--Xu already obtain good minimal models for projective klt pairs
of dimension at most four with nonzero Euler characteristic
\cite[Corollary~5.2]{LiuXu2025}; in the nef case this gives
semiampleness. Lazi\'c's work supplies a semiampleness criterion using
uniform multiplier-ideal comparisons, and an algebraic approximation
theorem for supercanonical currents, under their respective hypotheses
\cite[Theorems~A--B]{Lazic2024Metrics}.

In the compact K\"ahler setting, H\"oring--Lazi\'c--Lehn obtain
nonvanishing for non-uniruled $\Q$-factorial klt pairs with nef adjoint
of numerical dimension one and nonzero Euler characteristic
\cite[Theorem~A]{HoringLazicLehn2025}. Under these hypotheses in
dimension four, they obtain semiampleness when the minimal pullback
current has zero Lelong numbers
\cite[Corollary~B]{HoringLazicLehn2025}.

\Cref{sec:adjoint-metrics} proves the metric theorem, the arbitrary-data
localization estimate and the multiplier-ideal consequence.
\Cref{sec:interior-injectivity} proves the independent interior
injectivity theorem. \Cref{sec:euler-route} states the two published
fourfold criteria and applies them to the same minimal metric.
